% GENERATED by reporting/make_priors_report.py -- do not edit by hand.
% Source: metrics/outputs/{toy-temporal,gemma-2-2b*}/ -- frozen run results, not prose.
% Every caption is TODO on purpose: a caption is a claim about what to take away,
% and this project requires those to be written by a person.
% Regenerate:  python3 -m reporting.make_priors_report
% Requires:    \input{thresholds_macros.tex} in the preamble (defines \MinJoint etc.),
%              and packages booktabs, graphicx, amsmath.


% ------------------------------------------------------------------------
% gemma-2-2b layer 12
% ------------------------------------------------------------------------

% [gemma-layer]
\begin{table}[tbp]
  \centering
  \small
  \caption{TODO(caption, human-written). Fraction of variance explained at seven depths. An SAE reconstructs the layer it was fitted to best. The checkpoint's \texttt{conf.yaml} records \texttt{block\_id: 1} while the repository README lists layer 12.}
  \label{tab:pit-gemma-layer}
  \begin{tabular}{llr}
    \toprule
    activation & block output & FVE \\
    \midrule
    \texttt{hidden\_states[1]} & block 0 & -0.0207 \\
    \texttt{hidden\_states[2]} & block 1 & -0.0145 \\
    \texttt{hidden\_states[6]} & block 5 & -0.0245 \\
    \texttt{hidden\_states[12]} & block 11 & 0.7287 \\
    \texttt{hidden\_states[13]} & block 12 & \textbf{0.7932} \\
    \texttt{hidden\_states[18]} & block 17 & 0.5181 \\
    \texttt{hidden\_states[25]} & block 24 & 0.0876 \\
    \bottomrule
  \end{tabular}
  \\[2pt] \footnotesize Negative values mean the reconstruction is worse than predicting the mean. The peak is the output of block 12, matching the README.
\end{table}

% [gemma-quality]
\begin{table}[tbp]
  \centering
  \small
  \caption{TODO(caption, human-written). Released Priors in Time checkpoint at \texttt{gemma-2-2b}, output of block 12, on 200 documents of \texttt{pile-10k}. Measured with the authors' own forward pass.}
  \label{tab:pit-gemma-quality}
  \begin{tabular}{lr}
    \toprule
    quantity & value \\
    \midrule
    dictionary width & 9,216 \\
    tokens scored & 24,025 \\
    FVE & 0.7924 \\
    novel $L_0$ & 192.00 \footnotesize(TopK setting, $k=192$) \\
    dead features & 0 / 9,216 \\
    avg max cosine & 0.1143 \\
    \textbf{predictive share of reconstruction energy} & \textbf{0.904} \\
    \bottomrule
  \end{tabular}
  \\[2pt] \footnotesize Nine parts in ten of what reconstructs a token comes from its context. No hierarchy metric is reported: this method produces no feature edges, so the metric set has no matching input.
\end{table}

% ------------------------------------------------------------------------
% Toy model
% ------------------------------------------------------------------------

% [toy]
\begin{table}[tbp]
  \centering
  \small
  \caption{TODO(caption, human-written). Priors in Time trained on the temporal toy, three seeds per row. The null row uses \texttt{--persistence 0}, which removes the temporal correlation while leaving the per-token distribution unchanged.}
  \label{tab:pit-toy}
  \begin{tabular}{lcrr}
    \toprule
    condition & $n$ & FVE & $L_0$ \\
    \midrule
    temporal, original tree & 3 & 0.979\,$\pm$\,0.010 & 2.00\,$\pm$\,0.00 \\
    temporal, decorrelated tree & 3 & 0.926\,$\pm$\,0.009 & 2.00\,$\pm$\,0.00 \\
    i.i.d. (no events) & 3 & 0.965\,$\pm$\,0.007 & 2.00\,$\pm$\,0.00 \\
    \bottomrule
  \end{tabular}
  \\[2pt] \footnotesize $L_0$ is fixed by TopK at $k=2$ and is a setting rather than a measurement. The event-boundary measurements are in Table~\ref{tab:pit-boundary}.
\end{table}

% [toy-boundary]
\begin{table}[tbp]
  \centering
  \small
  \caption{TODO(caption, human-written). Novel-code magnitude and predictive-code similarity at planted parent-state changes, three seeds per row. Two predictions were stated before the run: novel activity higher at changes, and the difference absent in the null.}
  \label{tab:pit-boundary}
  \begin{tabular}{lrrrrr}
    \toprule
    condition & novel at & novel between & ratio & pred.\ cos at & pred.\ cos between \\
    \midrule
    temporal, original & 0.109 & 0.079 & 1.38 & 0.9309 & 0.9829 \\
    temporal, decorrelated & 0.140 & 0.061 & 2.31 & 0.7654 & 0.7987 \\
    \textbf{i.i.d. (no events)} & 0.118 & 0.073 & 1.62 & \textbf{0.4270} & 0.8954 \\
    \bottomrule
  \end{tabular}
  \\[2pt] \footnotesize The novel-code ratio exceeds one in the null as well, where the data carries no events. The predictive column separates the null from the other two. That measure was recorded alongside but not stated in advance.
\end{table}
